% Figure: Kantrowitz starting limit and the isentropic contraction
% limit for a supersonic inlet, as throat-to-capture area ratio versus
% flight Mach number. Curves are sampled and entered as coordinate
% lists; no in-axis trigonometry.
\begin{figure}[htbp]
\centering
\begin{tikzpicture}
\begin{axis}[
  width=12cm, height=8cm,
  xlabel={flight Mach number $M_\infty$},
  ylabel={throat-to-capture area ratio $A_t/A_c$},
  xmin=1, xmax=5, ymin=0, ymax=1.05,
  axis lines=left,
  grid=major, grid style={enggray!20},
  tick label style={font=\scriptsize},
  label style={font=\small},
  legend style={font=\scriptsize, at={(0.97,0.97)}, anchor=north east},
  legend cell align=left,
]
\addplot[engblue, very thick, mark=*, mark size=1.3pt] coordinates {
  (1.0,1.000) (1.5,0.768) (2.0,0.661) (2.5,0.620)
  (3.0,0.600) (3.5,0.589) (4.0,0.582) (5.0,0.574) };
\addlegendentry{Kantrowitz (self-start) limit}
\addplot[engred, very thick, mark=square*, mark size=1.3pt] coordinates {
  (1.0,1.000) (1.5,0.578) (2.0,0.373) (2.5,0.254)
  (3.0,0.181) (3.5,0.133) (4.0,0.100) (5.0,0.060) };
\addlegendentry{isentropic (max steady) contraction}
\node[engblack, font=\scriptsize, align=center, text width=2.6cm]
  at (axis cs:2.9,0.86) {self-starts at rest geometry};
\node[engblack, font=\scriptsize, align=center, text width=3.0cm]
  at (axis cs:3.4,0.42) {starts only with variable geometry or overspeed};
\node[engblack, font=\scriptsize, align=center, text width=3.0cm]
  at (axis cs:3.7,0.12) {cannot run steadily (unstartable)};
\end{axis}
\end{tikzpicture}
\caption[Kantrowitz and isentropic inlet limits]{The two contraction
limits that bound supersonic-inlet design, plotted as the
throat-to-capture area ratio against flight Mach number. A smaller area
ratio is a more aggressive internal contraction. The upper (Kantrowitz)
curve is the most aggressive fixed contraction that will self-start: an
inlet whose geometry sits above this curve swallows its starting normal
shock on its own. The lower (isentropic) curve is the most aggressive
contraction the flow will sustain once started, where the throat just
reaches $M=1$. Between the two curves the inlet runs well once started
but will not start from its running geometry, so it needs variable
geometry, a starting bleed, or an overspeed-then-throttle schedule.
Below the isentropic curve no steady supersonic solution exists and the
inlet is unstartable as built. The widening gap between the curves at
high Mach number is why high-Mach inlets carry the most elaborate
starting hardware.}
\label{fig:vol03ch13:kantrowitz}
\end{figure}
